
\documentclass{article}
\usepackage{geometry}
\geometry{margin=1.5cm, vmargin={0pt,1cm}}
\setlength{\topmargin}{-1cm}
\setlength{\headheight}{12.5pt}
\setlength{\paperheight}{29.7cm}
\setlength{\textheight}{25.3cm}

% useful packages.
% \usepackage[UTF8]{ctex}
\usepackage{fancyhdr}
\usepackage{layout}

% import common macros
% % % % % % % % % % % % % % % % % % % % % % % % % % % % % % % %
% Common Macros for LaTeX
% % % % % % % % % % % % % % % % % % % % % % % % % % % % % % % %

% Useful packages
\usepackage{amsfonts}
\usepackage{amsmath}
\usepackage{amssymb}
\usepackage{amsthm}
\usepackage{enumerate}
\usepackage{graphicx}
\usepackage{multicol}
\usepackage[ruled,linesnumbered]{algorithm2e}

% Math operators and common symbols
\newcommand{\dif}{\mathrm{d}}
\newcommand{\innerProd}[1]{\left\langle #1 \right\rangle}
\newcommand{\difFrac}[2]{\frac{\dif #1}{\dif #2}}
\newcommand{\pdfFrac}[2]{\frac{\partial #1}{\partial #2}}
\newcommand{\T}{\mathrm{T}}
\newcommand{\norm}[1]{\left\| #1 \right\|}
\newcommand{\abs}[1]{\left| #1 \right|}

% Vector and Matrix shortcuts
\newcommand{\vecTwo}[2]{
  \begin{bmatrix}
    #1 \\
    #2
\end{bmatrix}}
\newcommand{\vecThree}[3]{
  \begin{bmatrix}
    #1 \\
    #2 \\
    #3
\end{bmatrix}}
\newcommand{\vecTwoH}[2]{
  \begin{bmatrix}
    #1 & #2
  \end{bmatrix}
}
\newcommand{\vecThreeH}[3]{
  \begin{bmatrix}
    #1 & #2 & #3
  \end{bmatrix}
}

% Common fields
\newcommand{\R}{\mathbb{R}}
\newcommand{\C}{\mathbb{C}}
\newcommand{\Q}{\mathbb{Q}}
\newcommand{\Z}{\mathbb{Z}}
\newcommand{\N}{\mathbb{N}}

% Common Operators
\renewcommand{\Re}{\mathrm{Re}}
\renewcommand{\Im}{\mathrm{Im}}

% Theorem-like environments
\theoremstyle{definition}
\newtheorem{theorem}{Theorem}[section]
\newtheorem{lemma}[theorem]{Lemma}
\newtheorem{corollary}[theorem]{Corollary}
\newtheorem{definition}[theorem]{Definition}
\newtheorem{remark}[theorem]{Remark}
\newtheorem{exercise}{Exercise}[section]

% Support for individual Chinese characters using fontspec (for LuaLaTeX)
\usepackage{fontspec}
\newfontfamily\zhfont{SimSun} % Search system fonts by name
\newcommand{\zhstr}[1]{{\zhfont #1}}

% Solutions and proofs
\AtBeginDocument{
  \ifdefined\ctexset
  \newenvironment{solution}{
  \begin{proof}[解]}{
  \end{proof}}
  \else
  \newenvironment{solution}{
  \begin{proof}[Solution]}{
  \end{proof}}
  \fi
}

\usepackage{luacode}
% Utility to get environment variables (requires lualatex)
\newcommand{\getEnv}[2]{\directlua{tex.print(os.getenv("#1") or "#2")}}


% Name and uID
\newcommand{\name}{\getEnv{NAME}{NAME}}
\newcommand{\uid}{\getEnv{UID}{UID}}
\newcommand{\course}{Complex Analysis}
\newcommand{\hwNum}{2}

\begin{document}

\pagestyle{fancy}
\fancyhead{}
\lhead{\zhstr{\name} (\uid)}
\chead{\course\ \#\hwNum}
\rhead{\today}

\section{Exercise C.264}

\begin{lemma}[Lemma C.263]
  The signed volume $\delta$ is alternating, i.e.,
  \[
    \delta(v_1, \dots, v_i, \dots, v_j, \dots, v_n)
    = -\delta(v_1, \dots, v_j, \dots, v_i, \dots, v_n).
  \]
\end{lemma}

Prove Lemma C.263 using (SVP-2,3).

\begin{proof}
  We first prove the adjacent-swap case. Fix all slots except the
  $i$-th and $(i+1)$-th, and set $u := v_i + v_{i+1}$. Consider
  \[
    \delta(\dots, u, u, \dots).
  \]
  By (SVP-2), this equals $0$ because two entries are equal.
  By linearity (SVP-3) in each of the two slots, we expand:
  \[
    \begin{aligned}
      0
      &= \delta(\dots, v_i+v_{i+1}, v_i+v_{i+1}, \dots) \\
      &= \delta(\dots, v_i, v_i, \dots)
      + \delta(\dots, v_i, v_{i+1}, \dots) \\
      &\quad + \delta(\dots, v_{i+1}, v_i, \dots)
      + \delta(\dots, v_{i+1}, v_{i+1}, \dots).
    \end{aligned}
  \]
  Again by (SVP-2), the first and last terms are $0$. Hence
  \[
    \delta(\dots, v_i, v_{i+1}, \dots)
    = -\delta(\dots, v_{i+1}, v_i, \dots).
  \]

  Now let $i<j$. Swapping $v_i$ and $v_j$ can be done by adjacent
  transpositions:
  move $v_i$ right to position $j$ (total $j-i$ swaps), then move
  $v_j$ left to position $i$ (total $j-i-1$ swaps). So the total number
  of adjacent swaps is
  \[
    (j-i)+(j-i-1)=2(j-i)-1,
  \]
  which is odd. Each adjacent swap contributes a factor $-1$, therefore
  \[
    \delta(v_1,\dots,v_i,\dots,v_j,\dots,v_n)
    = -\delta(v_1,\dots,v_j,\dots,v_i,\dots,v_n).
  \]
  Thus $\delta$ is alternating.
\end{proof}

\section{Exercise C.344}

\begin{lemma}[Lemma C.343]
  For $f : E \to \R$ with $E \subset \R$, $x_0 \in E$, and $L \in \R$,
  the following two statements are equivalent:
  \begin{enumerate}
    \item $f$ is differentiable at $x_0$ and $f'(x_0)=L$;
    \item
      \[
        \lim_{x\to x_0,\,x\in E\setminus\{x_0\}}
        \frac{|f(x)-f(x_0)-L(x-x_0)|}{|x-x_0|}=0.
      \]
  \end{enumerate}
\end{lemma}

Prove Lemma C.343.

\begin{proof}
  (a) $\Rightarrow$ (b): If $f$ is differentiable at $x_0$ with derivative
  $L$, then by definition there exists a function $r(x)$ such that
  \[
    f(x)-f(x_0)=L(x-x_0)+r(x)(x-x_0),
    \qquad r(x)\to 0 \ (x\to x_0,\ x\in E\setminus\{x_0\}).
  \]
  Therefore, for $x\neq x_0$,
  \[
    \frac{|f(x)-f(x_0)-L(x-x_0)|}{|x-x_0|}=|r(x)|\to 0,
  \]
  which is exactly (b).

  (b) $\Rightarrow$ (a): Assume
  \[
    \lim_{x\to x_0,\,x\in E\setminus\{x_0\}}
    \frac{|f(x)-f(x_0)-L(x-x_0)|}{|x-x_0|}=0.
  \]
  Define for $x\neq x_0$,
  \[
    r(x):=\frac{f(x)-f(x_0)-L(x-x_0)}{x-x_0}.
  \]
  Then $|r(x)|\to 0$, hence $r(x)\to 0$, and we can rewrite
  \[
    f(x)-f(x_0)=L(x-x_0)+r(x)(x-x_0).
  \]
  This is exactly the differentiability condition at $x_0$ with
  derivative $L$. So (a) holds.

  Therefore (a) and (b) are equivalent.
\end{proof}

\section{Exercise C.353}

\begin{lemma}[Lemma C.352]
  Suppose a function $f : E \to \R^m$ with $E \subseteq \R^n$
  is differentiable at an interior point $x_0$ of $E$ with derivative
  $L_1$ and also differentiable at $x_0$ with derivative $L_2$.
  Then $L_1=L_2$.
\end{lemma}

Prove Lemma C.352.

\begin{proof}
  Since $f$ is differentiable at $x_0$ with derivative $L_1$, we have
  \[
    f(x_0+h)-f(x_0)=L_1h+r_1(h),
    \qquad \frac{|r_1(h)|}{|h|}\to 0 \ (h\to 0).
  \]
  Since $f$ is also differentiable at $x_0$ with derivative $L_2$, we have
  \[
    f(x_0+h)-f(x_0)=L_2h+r_2(h),
    \qquad \frac{|r_2(h)|}{|h|}\to 0 \ (h\to 0).
  \]
  Subtracting gives
  \[
    (L_1-L_2)h=r_2(h)-r_1(h).
  \]
  Hence
  \[
    \frac{|(L_1-L_2)h|}{|h|}
    \le
    \frac{|r_1(h)|}{|h|}+\frac{|r_2(h)|}{|h|}
    \to 0
    \quad (h\to 0,\ h\neq 0).
  \]

  Let $e_k$ be the $k$-th standard basis vector in $\R^n$, and set
  $h=t e_k$ with $t\neq 0$. Then
  \[
    \frac{|(L_1-L_2)(t e_k)|}{|t e_k|}
    = \frac{|t|\,|(L_1-L_2)e_k|}{|t|}
    = |(L_1-L_2)e_k|.
  \]
  Taking $t\to 0$ yields $|(L_1-L_2)e_k|=0$, so
  $(L_1-L_2)e_k=0$ for every $k=1,\dots,n$.

  Therefore $L_1-L_2$ vanishes on a basis of $\R^n$, hence
  $L_1-L_2=0$, i.e. $L_1=L_2$.
\end{proof}

\section{Exercise C.361}

Show that the existence of partial derivatives at $x_0$ does not imply
that the function is differentiable at $x_0$ by considering the
differentiability of the following function $f : \R^2 \to \R$ at
$(0,0)$:
\[
  f(x,y)=
  \begin{cases}
    \dfrac{x^3}{x^2+y^2}, & (x,y)\neq(0,0), \\
    0, & (x,y)=(0,0).
  \end{cases}
\]

\begin{solution}
  First compute the partial derivatives at $(0,0)$.
  For $\partial_x f(0,0)$,
  \[
    \partial_x f(0,0)
    = \lim_{h\to 0}\frac{f(h,0)-f(0,0)}{h}
    = \lim_{h\to 0}\frac{h-0}{h}=1.
  \]
  For $\partial_y f(0,0)$,
  \[
    \partial_y f(0,0)
    = \lim_{h\to 0}\frac{f(0,h)-f(0,0)}{h}
    = \lim_{h\to 0}\frac{0-0}{h}=0.
  \]
  So both partial derivatives exist at $(0,0)$.

  If $f$ were differentiable at $(0,0)$, the derivative would be
  \[
    L(x,y)=\partial_x f(0,0)x+\partial_y f(0,0)y=x.
  \]
  Now test the differentiability quotient along the path $y=x$:
  \[
    f(x,x)=\frac{x^3}{x^2+x^2}=\frac{x}{2},
    \qquad
    L(x,x)=x.
  \]
  Hence
  \[
    \frac{|f(x,x)-L(x,x)|}{\sqrt{x^2+x^2}}
    = \frac{|x/2-x|}{\sqrt{2}|x|}
    = \frac{1}{2\sqrt{2}}.
  \]
  This does not tend to $0$ as $x\to 0$. Therefore $f$ is not
  differentiable at $(0,0)$.

  We have shown that partial derivatives may exist at a point while
  differentiability fails at that point.
\end{solution}

\section{Exercise C.385}

Show that the unit sphere
\[
  S^2 := \{(x,y,z)\in\R^3 : x^2+y^2+z^2=1\}
\]
is a regular surface.

\begin{solution}
  Define
  \[
    F:\R^3\to\R, \qquad F(x,y,z):=x^2+y^2+z^2-1.
  \]
  Then
  \[
    S^2 = F^{-1}(0).
  \]
  Since $F$ is smooth, we compute
  \[
    \nabla F(x,y,z)=(2x,2y,2z).
  \]
  If $(x,y,z)\in S^2$, then $x^2+y^2+z^2=1$, so $(x,y,z)\neq(0,0,0)$, and
  therefore
  \[
    \nabla F(x,y,z)\neq 0 \quad \text{for all } (x,y,z)\in S^2.
  \]
  Hence $0$ is a regular value of $F$. By the regular level-set theorem
  (equivalently, the implicit function theorem), $F^{-1}(0)$ is a regular
  surface of $\R^3$.

  So $S^2$ is a regular surface.
\end{solution}

\section{Exercise C.388}

\begin{lemma}[Lemma C.387]
  Let $x : U\subset\R^2 \to S$ be a parametrization of a regular surface
  $S$ and let $q\in U$. The vector space
  \[
    dx_q(\R^2)\subset\R^3
  \]
  coincides with the set of tangent vectors to $S$ at $x(q)$.
\end{lemma}

Prove Lemma C.387.

\begin{proof}
  Let $p:=x(q)$.

  We first show that every tangent vector at $p$ belongs to $dx_q(\R^2)$.
  Let $w$ be a tangent vector to $S$ at $p$. By definition, there exists
  a differentiable curve $\alpha:(-\varepsilon,\varepsilon)\to S$ such that
  \[
    \alpha(0)=p, \qquad \alpha'(0)=w.
  \]
  Since $x:U\to V\cap S$ is a local parametrization and a homeomorphism,
  for $t$ near $0$ we may write
  \[
    \alpha(t)=x(\beta(t)), \qquad
    \beta(t):=x^{-1}(\alpha(t))\in U,
  \]
  with $\beta(0)=q$. By the chain rule,
  \[
    w=\alpha'(0)=d(x\circ\beta)_0=dx_q\big(\beta'(0)\big)\in dx_q(\R^2).
  \]
  Hence the set of tangent vectors is contained in $dx_q(\R^2)$.

  Conversely, let $u\in dx_q(\R^2)$. Then there exists $v\in\R^2$ such that
  $u=dx_q(v)$. Define
  \[
    \beta(t):=q+tv
  \]
  for $|t|$ small enough so that $\beta(t)\in U$, and set
  \[
    \alpha(t):=x(\beta(t))=x(q+tv)\in S.
  \]
  Then $\alpha$ is a differentiable curve in $S$ with $\alpha(0)=x(q)=p$, and
  \[
    \alpha'(0)=dx_q(v)=u.
  \]
  So $u$ is a tangent vector to $S$ at $p$.

  Therefore $dx_q(\R^2)$ coincides with the set of tangent vectors to $S$
  at $x(q)$.
\end{proof}

\section{Exercise D.6}

Let $X$ be the set of all bounded and unbounded sequences of complex
numbers. Show that the function $d$ given by
\[
  \forall x=(\xi_j),\ \forall y=(\eta_j),\
  d(x,y)=\sum_{j=1}^{\infty}\frac{1}{2^j}\frac{|\xi_j-\eta_j|}{1+|\xi_j-\eta_j|}
\]
is a metric on $X$.

\begin{solution}
  Let
  \[
    \phi(t):=\frac{t}{1+t},\qquad t\ge 0.
  \]
  Then
  \[
    d(x,y)=\sum_{j=1}^{\infty}2^{-j}\phi\big(|\xi_j-\eta_j|\big).
  \]
  Since $0\le \phi(t)<1$, each term is between $0$ and $2^{-j}$, so the
  series converges absolutely by comparison with $\sum 2^{-j}$.

  \textbf{(1) Non-negativity.} Each summand is nonnegative, hence
  $d(x,y)\ge 0$.

  \textbf{(2) Identity of indiscernibles.}
  If $x=y$, then every $|\xi_j-\eta_j|=0$, so $d(x,y)=0$.
  Conversely, if $d(x,y)=0$, then a sum of nonnegative terms is zero, so each
  term is zero. Thus
  \[
    \phi\big(|\xi_j-\eta_j|\big)=0 \Rightarrow |\xi_j-\eta_j|=0
    \Rightarrow \xi_j=\eta_j
  \]
  for all $j$, hence $x=y$.

  \textbf{(3) Symmetry.} Since $|\xi_j-\eta_j|=|\eta_j-\xi_j|$, we get
  $d(x,y)=d(y,x)$.

  \textbf{(4) Triangle inequality.}
  For $x=(\xi_j), y=(\eta_j), z=(\zeta_j)$, by the usual triangle inequality,
  \[
    |\xi_j-\zeta_j|\le |\xi_j-\eta_j|+|\eta_j-\zeta_j|.
  \]
  We claim that for $a,b\ge 0$,
  \[
    \phi(a+b)\le \phi(a)+\phi(b).
  \]
  Indeed,
  \[
    \phi(a)+\phi(b)-\phi(a+b)
    = \frac{a}{1+a}+\frac{b}{1+b}-\frac{a+b}{1+a+b}
    = \frac{ab(a+b+2)}{(1+a)(1+b)(1+a+b)}\ge 0.
  \]
  Hence
  \[
    \phi(|\xi_j-\zeta_j|)
    \le \phi(|\xi_j-\eta_j|+|\eta_j-\zeta_j|)
    \le \phi(|\xi_j-\eta_j|)+\phi(|\eta_j-\zeta_j|).
  \]
  Multiply by $2^{-j}$ and sum over $j$:
  \[
    d(x,z)\le d(x,y)+d(y,z).
  \]

  Therefore $d$ satisfies all metric axioms, so $d$ is a metric on $X$.
\end{solution}

\section{Exercise D.18}

The completeness depends on the metric. For the metric
\[
  d_1(x,y)=\int_a^b |x(t)-y(t)|\,dt,
\]
show that the metric space $(C[a,b], d_1)$ is not complete.

\begin{solution}
  Let
  \[
    c:=\frac{a+b}{2}, \qquad \varepsilon_n:=\frac{b-a}{4n}.
  \]
  Define $f_n\in C[a,b]$ by
  \[
    f_n(t):=
    \begin{cases}
      0, & a\le t\le c-\varepsilon_n, \\
      \dfrac{t-(c-\varepsilon_n)}{2\varepsilon_n},
      & c-\varepsilon_n\le t\le c+\varepsilon_n, \\
      1, & c+\varepsilon_n\le t\le b.
    \end{cases}
  \]
  This is a continuous piecewise linear function.

  Define the step function
  \[
    H(t):=
    \begin{cases}
      0, & t<c, \\
      1, & t\ge c.
    \end{cases}
  \]
  Then $H$ is not continuous at $c$.

  We estimate $d_1(f_n,H)$. Outside $[c-\varepsilon_n,c+\varepsilon_n]$,
  we have $f_n=H$. So
  \[
    d_1(f_n,H)
    =\int_{c-\varepsilon_n}^{c+\varepsilon_n}|f_n(t)-H(t)|\,dt
    =\frac{\varepsilon_n}{2}
    =\frac{b-a}{8n}\to 0.
  \]
  Hence for $m,n$,
  \[
    d_1(f_n,f_m)\le d_1(f_n,H)+d_1(f_m,H)\to 0,
  \]
  so $(f_n)$ is Cauchy in $(C[a,b],d_1)$.

  We now show $(f_n)$ has no limit in $C[a,b]$. Assume by contradiction
  that there exists $g\in C[a,b]$ with $d_1(f_n,g)\to 0$.

  If $g(c)\le \tfrac12$, continuity gives some $\delta>0$ such that
  $g(t)\le \tfrac34$ for $|t-c|<\delta$. For $n$ large enough,
  $\varepsilon_n<\delta/2$, and on
  \[
    I:=[c+\delta/2,\,c+\delta]\subset[c+\varepsilon_n,b],
  \]
  we have $f_n=1$. Thus
  \[
    |f_n-g|=|1-g|\ge \tfrac14 \quad \text{on } I,
  \]
  so
  \[
    d_1(f_n,g)\ge \int_I |f_n-g|\,dt\ge \frac{\delta}{8},
  \]
  contradicting $d_1(f_n,g)\to 0$.

  If $g(c)>\tfrac12$, continuity gives some $\delta>0$ such that
  $g(t)\ge \tfrac14$ for $|t-c|<\delta$. For $n$ large enough,
  $\varepsilon_n<\delta/2$, and on
  \[
    J:=[c-\delta,\,c-\delta/2]\subset[a,c-\varepsilon_n],
  \]
  we have $f_n=0$. Hence
  \[
    |f_n-g|=|g|\ge \tfrac14 \quad \text{on } J,
  \]
  and therefore
  \[
    d_1(f_n,g)\ge \int_J |f_n-g|\,dt\ge \frac{\delta}{8},
  \]
  again a contradiction.

  So no continuous $g$ can be the $d_1$-limit of $(f_n)$. Therefore
  $(C[a,b],d_1)$ is not complete.
\end{solution}

\end{document}
